% GENERATED FILE. Edit canonical lessons, then run npm run generate:books.
% canonical-source-hash: fnv1a64:840cc974f1c53efd
% canonical-lessons: TA-C69-vehicle, TA-C69-bus, TA-W25-read-vandi, TA-C69-train, TA-C69-road, TA-C69-station

\chapter{Getting There}
\label{ch:ta-getting-there}
\hlchaptermodality{69}

\begin{chapteropening}
\textbf{By the end of this chapter:} \emph{I can name a vehicle, a bus, a train, a road and a station in Tamil.}

Complete TA-C69-station and say all five words of the chapter in order.
\end{chapteropening}

\section[vaṇḍi]{\ta{வண்டி} (vaṇḍi) — a vehicle, a cart}

\label{lesson:TA-C69-vehicle}



\begin{quote}
Before the new one: what did \emph{pai} mean?
\end{quote}

\subsection*{You'll want to know: \ta{வண்டி}}
\textbf{\ta{வண்டி}} (\emph{vaṇḍi}) — \textquotedblleft{}a vehicle, a cart\textquotedblright{}.

The general word for a thing with wheels. A bullock cart is a \ta{வண்டி}; so is a bicycle, and so is a car. Tamil starts here and adds a word in front when it needs to be exact.

In the middle sits \textbf{\ta{ண்ட}} — the curled-back \ta{ண} closed by a \ta{புள்ளி} and leaning on \ta{ட}. That pair is one of the most common shapes in the language, and you will now start seeing it everywhere.

\begin{grammarlens}[title={What you've built}]
The word the rest of this chapter is built on.
\end{grammarlens}

\subsection*{Your turn}
\begin{itemize}
\raggedright
  \item \emph{Say it:} \emph{vaṇḍi}
  \item \emph{Say it:} it once more, curling the tongue back for the \ta{ண}
  \item \emph{Recall:} say \emph{vāṅgu}, then say \emph{pai}, then say \emph{vaṇḍi}
\end{itemize}

\begin{tcolorbox}[breakable,colback=teal!4,colframe=teal!35!black,title={Before you move on}]
What does \ta{வண்டி} mean? (\textquotedblleft{}A vehicle, a cart\textquotedblright{}.)
\end{tcolorbox}
\section[pērundu]{\ta{பேருந்து} (pērundu) — a bus}

\label{lesson:TA-C69-bus}



\begin{quote}
Before the new one: what did \emph{vaṇḍi} mean?
\end{quote}

\subsection*{You'll want to know: \ta{பேருந்து}}
\textbf{\ta{பேருந்து}} (\emph{pērundu}) — \textquotedblleft{}a bus\textquotedblright{}.

A made word, and you can take it apart. \textbf{\ta{பேர்}-} is the big-word \ta{பெரு}, and \textbf{\ta{உந்து}} is \textquotedblleft{}to push along\textquotedblright{}. A big pushed thing.

This is the printed and announced word: it is what stands on the board at the stop and in the newspaper. In the street you will also hear the English \textquotedblleft{}bus\textquotedblright{} said plainly, and both are understood — but only one of them is written on the board you have to read.

\begin{grammarlens}[title={What you've built}]
Two, and one of them you can take apart.
\end{grammarlens}

\subsection*{Your turn}
\begin{itemize}
\raggedright
  \item \emph{Say it:} \emph{pērundu}
  \item \emph{Say it:} the two pieces, then the whole word
  \item \emph{Recall:} say \emph{pai}, then say \emph{vaṇḍi}, then say \emph{pērundu}
\end{itemize}

\begin{tcolorbox}[breakable,colback=teal!4,colframe=teal!35!black,title={Before you move on}]
What are the two pieces inside \ta{பேருந்து}? (\textquotedblleft{}\textbf{Big}\textquotedblright{}, and \textquotedblleft{}\textbf{push along}\textquotedblright{}.)
\end{tcolorbox}
\section[vaṇḍi]{\ta{வண்டி} — a dot joining two letters}

\label{lesson:TA-W25-read-vandi}



\begin{quote}
Before the page: you can say \emph{vaṇḍi} and \emph{pērundu}. Now look at the first one on the page.
\end{quote}

\subsection*{Script you'll notice: \ta{வண்டி} — a dot joining two letters}
\textbf{\ta{வண்டி}} brings no new shape either. What it brings is the \ta{புள்ளி} doing its job \textbf{between} two different letters, where \ta{சட்டை} showed it between two of the same one.

\noindent
\begin{tabularx}{\linewidth}{@{}>{\raggedright\arraybackslash}X>{\raggedright\arraybackslash}X>{\raggedright\arraybackslash}X@{}}
\toprule
\textbf{piece} & \textbf{} & \textbf{} \\
\midrule
\textbf{\ta{வ}} & \emph{va} & the letter with its built-in \emph{a} \\
\textbf{\ta{ண்}} & \ta{ண} with a \textbf{\ta{புள்ளி}} & closed — nothing said after it \\
\textbf{\ta{டி}} & \emph{ḍi} & \ta{ட} with the \textbf{\ta{ி}} sign \\
\bottomrule
\end{tabularx}

\begin{quote}
\textbf{\ta{வ} · \ta{ண்} · \ta{டி} $\to$ \ta{வண்டி}}
\end{quote}

Read the middle two as one movement rather than two beats. That is what the \ta{புள்ளி} is asking for: \ta{ண்} has no vowel of its own left, so it has nowhere to go except onto the letter after it.

\subsection*{Your turn}
\begin{itemize}
\raggedright
  \item \emph{Read it:} \textbf{\ta{வ}} — then \textbf{\ta{ண்}} — then \textbf{\ta{டி}}
  \item \emph{Point to:} the \ta{புள்ளி}, and say which letter it has emptied
  \item \emph{Read it:} \textbf{\ta{சட்டை}}, then \textbf{\ta{வண்டி}} — the same dot, two different jobs
  \item \emph{Recall:} say \emph{vaṇḍi}, then say \emph{pērundu}, then read \textbf{\ta{வண்டி}}
\end{itemize}

\begin{tcolorbox}[breakable,colback=teal!4,colframe=teal!35!black,title={Before you move on}]
Read \textbf{\ta{வண்டி}}. What has the \ta{புள்ளி} done to the \ta{ண}? (Taken its vowel away, so it leans on the \ta{ட}.)
\end{tcolorbox}
\section[rayil]{\ta{ரயில்} (rayil) — a train}

\label{lesson:TA-C69-train}



\begin{quote}
Before the new one: read \textbf{\ta{வண்டி}}, and say what it means.
\end{quote}

\subsection*{You'll want to know: \ta{ரயில்}}
\textbf{\ta{ரயில்}} (\emph{rayil}) — \textquotedblleft{}a train\textquotedblright{}.

A borrowed word, and it wears the borrowing openly: English \textquotedblleft{}rail\textquotedblright{}, written straight out in Tamil letters.

There is a clue in the very first letter. Tamil words do not normally \textbf{begin} with \ta{ர}, so a word that starts that way has almost always come in from somewhere else. You now have a small test you can run on a word before anyone tells you where it is from.

\begin{grammarlens}[title={What you've built}]
Three, and one of them is a visitor.
\end{grammarlens}

\subsection*{Your turn}
\begin{itemize}
\raggedright
  \item \emph{Say it:} \emph{rayil}
  \item \emph{Say it:} \emph{vaṇḍi}, then \emph{rayil} — and say which one is Tamil's own
  \item \emph{Recall:} say \emph{pērundu}, then read \textbf{\ta{வண்டி}}, then say \emph{rayil}
\end{itemize}

\begin{tcolorbox}[breakable,colback=teal!4,colframe=teal!35!black,title={Before you move on}]
What does the \ta{ர} at the front tell you about \ta{ரயில்}? (That it is a \textbf{borrowed} word.)
\end{tcolorbox}
\section[sālai]{\ta{சாலை} (sālai) — a road}

\label{lesson:TA-C69-road}



\begin{quote}
Before the new one: what did \emph{rayil} mean?
\end{quote}

\subsection*{You'll want to know: \ta{சாலை}}
\textbf{\ta{சாலை}} (\emph{sālai}) — \textquotedblleft{}a road\textquotedblright{}.

A made road, wide enough for a \ta{வண்டி} to go down.

Set it beside \ta{பாதை} (\emph{pātai}), the path you already have. A \ta{பாதை} is walked; a \ta{சாலை} is driven. The difference is not size so much as what it was made for.

Both end in the same \textbf{\ta{ை}}, which by now you have read four times.

\begin{grammarlens}[title={What you've built}]
Four, and the vehicles have somewhere to run.
\end{grammarlens}

\subsection*{Your turn}
\begin{itemize}
\raggedright
  \item \emph{Say it:} \emph{sālai}
  \item \emph{Say it:} \emph{pātai}, then \emph{sālai} — and say which one you drive down
  \item \emph{Recall:} read \textbf{\ta{வண்டி}}, then say \emph{rayil}, then say \emph{sālai}
\end{itemize}

\begin{tcolorbox}[breakable,colback=teal!4,colframe=teal!35!black,title={Before you move on}]
What is the difference between a \ta{பாதை} and a \ta{சாலை}? (One is \textbf{walked}; one is \textbf{driven}.)
\end{tcolorbox}
\section[nilaiyam]{\ta{நிலையம்} (nilaiyam) — a station}

\label{lesson:TA-C69-station}



\begin{quote}
Before the new one: what did \emph{sālai} mean, and what did \emph{rayil} mean?
\end{quote}

\subsection*{You'll want to know: \ta{நிலையம்}}
\textbf{\ta{நிலையம்}} (\emph{nilaiyam}) — \textquotedblleft{}a station\textquotedblright{}.

The place where the vehicle stands still.

Inside it is \ta{நிலை} (\emph{nilai}) — \textquotedblleft{}standing, state\textquotedblright{} — and you have met that piece before without being told. It is sitting inside \ta{வானிலை} (\emph{vāṉilai}), the weather, which is literally the sky's \emph{state}. One small piece, two words that could hardly be further apart.

That is worth more than the word itself: Tamil builds long words out of short ones you already own, and from here on you can start looking for the pieces.

\begin{grammarlens}[title={What you've built}]
Five, and the chapter runs end to end.
\end{grammarlens}

\subsection*{Your turn}
\begin{itemize}
\raggedright
  \item \emph{Say it:} \emph{nilaiyam}
  \item \emph{Say it:} \emph{vāṉilai}, then \emph{nilaiyam}, and name the piece they share
  \item \emph{Say it:} all five — \emph{vaṇḍi}, \emph{pērundu}, \emph{rayil}, \emph{sālai}, \emph{nilaiyam}
  \item \emph{Recall:} say \emph{rayil}, then say \emph{sālai}, then say \emph{nilaiyam}
\end{itemize}

\begin{tcolorbox}[breakable,colback=teal!4,colframe=teal!35!black,title={Before you move on}]
Which piece do \ta{வானிலை} and \ta{நிலையம்} share, and what does it mean? (\textbf{\ta{நிலை}} — \textquotedblleft{}standing, state\textquotedblright{}.)
\end{tcolorbox}
